% Datos de portada, cabeceras y metadatos del Tema 4
\newcommand{\PortadaAsignatura}{Electromagnetismo I}
\newcommand{\PortadaTema}{Tema 4}
\newcommand{\PortadaTitulo}{Materiales\\[2pt]dieléctricos}
\newcommand{\PortadaSubtitulo}{Polarización y carga ligada, vector desplazamiento,\\
  medios lineales, condiciones de frontera y cavidades}
\newcommand{\PortadaImagen}{\ifoscuro portada-oscuro\else portada-claro\fi}
\newcommand{\PortadaFormulas}{\rhoP=-\divg\vP,\qquad \sigmaP=\vP\cdot\nn,\qquad
  \divg\vD=\rhof,\qquad \vD=\eo\vE+\vP}
\newcommand{\PortadaAutor}{Daniel Martín Guerrero}
\newcommand{\PortadaProfesor}{Luis López Díaz}
\newcommand{\PortadaTipo}{Apuntes formalizados}
\newcommand{\PortadaCurso}{Salamanca, curso 2026/2027}
\renewcommand{\cabizq}{Tema 4 $\cdot$ Materiales dieléctricos}
\renewcommand{\cabder}{Electromagnetismo I}
\hypersetup{pdftitle={Tema 4. Materiales dieléctricos},
  pdfauthor={Daniel Martín Guerrero},
  pdfsubject={Electromagnetismo I. Apuntes formalizados con seminarios prácticos},
  pdfkeywords={dieléctrico, polarización, vector desplazamiento, condiciones de frontera, cavidades}}
